\section{Configuração do Ambiente de Desenvolvimento}

Para executar o projeto LibrasVision, é necessário configurar o ambiente com as seguintes dependências:

\begin{lstlisting}[language=bash]
# Criar ambiente virtual Python
python -m venv venv

# Ativar ambiente (Linux/Mac)
source venv/bin/activate

# Ou (Windows)
venv\Scripts\activate

# Instalar dependências
pip install opencv-python mediapipe tensorflow numpy pandas scikit-learn matplotlib
\end{lstlisting}

\section{Arquitetura do Dataset}

O dataset foi organizado da seguinte forma:

\begin{verbatim}
dataset/
├── sinais/
│   ├── biblioteca/
│   │   ├── usuario1_01.mp4
│   │   ├── usuario1_02.mp4
│   │   └── ...
│   ├── prova/
│   │   ├── usuario1_01.mp4
│   │   └── ...
│   └── secretaria/
│       └── ...
├── landmarks/
│   ├── biblioteca.pkl
│   ├── prova.pkl
│   └── secretaria.pkl
└── labels.csv
\end{verbatim}

\section{Hiperparâmetros do Modelo LSTM}

Os hiperparâmetros utilizados no treinamento do modelo LSTM foram os seguintes:

\begin{table}[htb]
    \centering
    \caption{Hiperparâmetros do modelo LSTM}
    \label{tab:hiperparametros}
    \begin{tabular}{|l|c|}
        \hline
        \textbf{Parâmetro} & \textbf{Valor} \\
        \hline
        Número de unidades LSTM (1ª camada) & 64 \\
        Número de unidades LSTM (2ª camada) & 32 \\
        Taxa de aprendizado (learning rate) & 0.001 \\
        Número de épocas & 100 \\
        Tamanho do batch & 32 \\
        Função de ativação & ReLU \\
        Função de perda & Categorical Crossentropy \\
        Otimizador & Adam \\
        \hline
    \end{tabular}
\end{table}

\section{Matriz de Confusão do Modelo}

A matriz de confusão resultante da validação do modelo apresentou os seguintes resultados:

\begin{table}[htb]
    \centering
    \caption{Matriz de confusão - Reconhecimento de sinais}
    \label{tab:matriz_confusao}
    \begin{tabular}{|c|c|c|c|}
        \hline
        \textbf{Sinal Real} & \textbf{Biblioteca} & \textbf{Prova} & \textbf{Secretaria} \\
        \hline
        Biblioteca & 47 & 2 & 1 \\
        Prova & 1 & 49 & 0 \\
        Secretaria & 0 & 1 & 50 \\
        \hline
    \end{tabular}
\end{table}

\section{Exemplo de Landmarks Extraídos}

A estrutura dos landmarks extraídos pelo MediaPipe Holistic segue o seguinte padrão:

\begin{lstlisting}[language=Python]
# Estrutura dos landmarks
landmarks = {
    'pose': [  # 33 pontos do corpo
        {'x': 0.5, 'y': 0.3, 'z': -0.1, 'visibility': 0.99},
        # ... mais 32 pontos
    ],
    'left_hand': [  # 21 pontos da mão esquerda
        {'x': 0.4, 'y': 0.35, 'z': 0.05, 'visibility': 0.95},
        # ... mais 20 pontos
    ],
    'right_hand': [  # 21 pontos da mão direita
        {'x': 0.6, 'y': 0.35, 'z': 0.05, 'visibility': 0.95},
        # ... mais 20 pontos
    ],
    'face': [  # 468 pontos do rosto
        {'x': 0.5, 'y': 0.2, 'z': -0.1, 'visibility': 0.98},
        # ... mais 467 pontos
    ]
}
\end{lstlisting}

\section{Métricas de Desempenho}

Além da acurácia, foram calculadas as seguintes métricas de desempenho:

\begin{table}[htb]
    \centering
    \caption{Métricas de desempenho por classe}
    \label{tab:metricas_desempenho}
    \begin{tabular}{|l|c|c|c|}
        \hline
        \textbf{Sinal} & \textbf{Precisão} & \textbf{Recall} & \textbf{F1-Score} \\
        \hline
        Biblioteca & 0.98 & 0.94 & 0.96 \\
        Prova & 0.96 & 0.98 & 0.97 \\
        Secretaria & 0.98 & 0.98 & 0.98 \\
        \hline
        Média & 0.97 & 0.97 & 0.97 \\
        \hline
    \end{tabular}
\end{table}

\section{Considerações sobre Generalização}

A análise de generalização do modelo revelou que:

\begin{itemize}
    \item O modelo apresenta bom desempenho em variações moderadas de iluminação;
    \item Usuários com características físicas semelhantes aos do dataset obtiveram acurácia acima de 90\%;
    \item Ângulos de câmera não tão radicais quanto aos do treinamento mantêm desempenho superior a 85\%;
    \item A latência média de resposta ficou em aproximadamente 80-100ms, adequada para aplicações interativas.
\end{itemize}